\section{Experiments}
\label{sec:exp}

We benchmark \name against recent state-of-the-art methods.
After inspecting qualitative differences in capability in \cref{sec:exp:qualitative}, we proceed with evaluating \task performance in \cref{sec:exp:4d}, focusing the comparison on the best methods from \cref{sec:related:tracking}.
We then continue with pure reconstruction tasks in \cref{sec:exp:3d}, comparing to the recent SOTA models mentioned in \cref{sec:related:recon}.
We conclude with a set of key ablations in \cref{sec:exp:ablations}.

\para{Training setup.}
For the encoder, we use the ViT-g model variant with 40 layers on a spatio-temporal patch size of $2{\times}16{\times}16$.
This encoder and our 8-layer cross-attention decoder contain 1\,B and 144\,M parameters, respectively.
Our training mixture contains public and internal datasets including
BlendedMVS~\cite{yao2020blendedmvs},
Co3Dv2~\cite{reizenstein2021common},
Dynamic Replica~\cite{karaev2023dynamicstereo},
Kubric~\cite{movi-e},
MVS-Synth~\cite{DeepMVS},
PointOdyssey~\cite{PointOdyssey}, ScanNet++~\cite{yeshwanthliu2023scannetpp},
ScanNet~\cite{dai2017scannet},
Tartanair~\cite{tartanair2020iros},
VirtualKitti~\cite{virtual_kitti}, and
Waymo Open~\cite{Sun_2020_CVPR}.
We train on 48-frame clips at $256{\times}256$ resolution, decoding 2048 random queries that are oversampled on specific regions for more efficient training
, see the appendix for details.
The model is trained for 500\,k steps using the AdamW optimizer~\cite{adam, adamw} with a local batch size of 1 across 64 TPU chips, taking a total of just over 2 days to complete.

\input{figs/fps}
\input{tabs/3d_tracking}


\subsection{Qualitative Analysis}
\label{sec:exp:qualitative}

We begin our evaluations by inspecting video reconstruction quality in \cref{fig:point_cloud_vis}.
Dynamic objects (\eg, swan and train) reveal qualitative differences, highlighting performance gaps between different methods.
Pure reconstruction methods (see \cref{sec:related:recon}), exemplified by the state-of-the-art models MegaSaM and $\pi^3$, fail to understand dynamics in the scenes, producing either repeated images of the entities as they move through the scene, or failing to reconstruct them altogether.
Conversely, tracking-based models successfully reconstruct and track dynamics.
However, they commonly only track points from one frame, leaving gaps in the reconstruction in all areas that are occluded in the first frame as seen for SpatialTrackerV2~\cite{xiao2025spatialtrackerv2} here.
Meanwhile, \name is able to track all dynamic pixels of the video in a unified reference frame.
We further evaluate our method on diverse in-the-wild videos in \cref{fig:ours_vis}, demonstrating its performance to both static and dynamic scenes.


\subsection{\task}
\label{sec:exp:4d}

\input{tabs/video_depth}

We now turn to the quantitative evaluation of \task, comparing \name to other methods that can estimate dynamic correspondences.
We evaluate on TAPVid-3D~\cite{koppula2024tapvid} which is comprised of three subsets of real-world videos in challenging settings.

We first evaluate 3D tracking in the local camera frame using the standard protocol from \citet{koppula2024tapvid}, which measures a model's ability to predict and track the 3D position of every observed point in the \emph{local} camera reference frames.
We report standard 3D tracking metrics \emph{Average percent of points within delta error} $\text{APD}_{3D}$, \emph{Occlusion Accuracy} (OA) and \emph{3D Average Jaccard} (AJ).
As shown in \cref{tab:3d_track} (left), \name achieves state-of-the-art results against competing methods, both with and without known ground-truth intrinsics.

\cref{tab:3d_track} (right) shows results for \emph{world coordinate 3D tracking}.
This task measures a model's ability to predict tracks within a single, consistent world coordinate system, thereby also measuring the models' ability to implicitly change reference frames.
PStudio is omitted here as it has no camera motion, rendering this task effectively identical to the preceding one.
Since there is no concept of occlusion for this task, we report $\text{APD}_{3D}$ as well as the L1 distance between the predicted and ground-truth tracks which provides a comprehensive, global measure of deviation.
We observe that our model excels at this task as well, exhibiting strong improvements across metrics.

We finally compare efficiency in
\cref{tab:fps}, by counting the number of tracks each model can produce for a given target frame-per-second rate (FPS).
The measurements show that our model is 18-300$\times$ faster than prior methods.

\input{tabs/camera_pose}

\subsection{3D Reconstruction}
\label{sec:exp:3d}

We now evaluate the 3D reconstruction capabilities of our model.
The most holistic task here is 3D point cloud reconstruction, which requires predicting all observed pixels in the same world coordinate system.

We evaluate on two standard benchmarks: MPI Sintel~\cite{bulter2012sintel} and ScanNet~\cite{scannet}, which consist of dynamic and static scenes, respectively.
For evaluation, we first align the predicted and ground-truth point clouds via mean-shifting following \citet{li2025megasam} before reporting the mean L1 distance.
As shown in \cref{tab:depth} (left), our model outperforms all recent state-of-the-art models across datasets.



\para{Depth estimation.}
We evaluate depth estimation performance across four datasets: Sintel~\cite{bulter2012sintel},
ScanNet~\cite{scannet},
KITTI~\cite{kitti2013}, and
Bonn~\cite{Bonn2019}.
Following prior works~\cite{li2025megasam, wang2025cut3r}, we align the predicted video depth with the ground-truth depth before computing metrics.
Specifically, we adopt two alignment settings:
\emph{scale-only} (S) and \emph{scale-and-shift} (SS) alignment.
Both settings determine a single global scale factor between the predicted and ground-truth depths across the sequence, while the latter introduces an additional 3D translation term, following an affine-invariant setup.
After alignment, we compute the \emph{Absolute Relative Error} (AbsRel)~\cite{AbsRelErr} between the predicted and ground-truth depths.




We show results in \cref{tab:depth} (right).
\name achieves top-tier performance across all benchmarks, with particularly strong results on Sintel which is the most challenging dataset due to its highly dynamic scenes.



\para{Camera pose estimation.}
We evaluate camera pose estimation performance on the same datasets.
Following prior work~\cite{wang2025cut3r, jiang2025geo4d}, we evaluate on the $14$-sequence subset of Sintel which provides full rendering with motion blur and atmospheric effects to measure in-the-wild performance.
We report the \emph{Absolute Translation Error} (ATE), \emph{Relative Translation Error} (RPE trans), and \emph{Relative Rotation Error} (RPE rot) after Sim(3) alignment with the ground truth as in \cite{zhang2024monst3r, wang2025cut3r, wang2025pi}, as well as the \emph{Pose AUC@$30$} as in \cite{wang2025vggt}.
\cref{tab:camera} presents the results, where our model strongly outperforms competing baselines on both indoor scenes (ScanNet and Re10K) and simulated outdoor scenes (Sintel).




\cref{fig:speed} visualizes efficiency and accuracy for camera pose estimation across several recent state-of-the-art methods.
To ensure a fair comparison, we remove the decoding heads of the baselines which are unrelated to camera estimation, making them faster.
Despite this, we find that \name surpasses all existing methods in both accuracy and efficiency, notably outperforming MegaSaM in throughput by two orders of magnitude.

\subsection{Ablation Studies}
\label{sec:exp:ablations}

We present a series of ablation studies to validate our design choices and hyperparameters.
Unless otherwise noted, we adopt a ViT-L backbone with $6$ decoder layers as the default setting.
The model is trained for 300\,k iterations with a batch size 16.
All results are reported on the challenging Sintel dataset for both video depth and camera pose.


\input{figs/ablation_source_patch}
\input{tabs/ablation_source_patch}


\para{Local RGB patch.}
As described in \cref{sec:method:architecture}, we introduce an additional embedding that encodes the local appearance around the query's source location.
This additional embedding yields dramatic performance improvements as shown in \cref{tab:abla_source_patch} and \cref{fig:ablation_source_patch}. These improvements can be attributed to two factors: (i) the local appearance information helps the query establish more reliable correspondences with the encoded spatiotemporal features; and (ii) these patches provide low-level cues that help to segment objects from their surroundings, leading to fine-grained predictions, as exemplified by the sharper depth estimates. We note that while most related works have specialized modules to preserve low level details, such as DPT~\cite{Ranftl2021} with its skip connections between encoder and decoder layers, our approach is much simpler.
In \cref{sec:app:high-res} of the appendix, we further highlight how the patches unlock subpixel precision.


\para{Encoder size.}
We now examine how performance scales with the size of our encoder.
\cref{tab:abla_model_size} summarizes our findings for ViT-B (90\,M) to ViT-g (1\,B).
As the number of trainable parameters increases, we observe a significant performance improvement, particularly in depth estimation and RPE-R in camera pose estimation.

\input{tabs/ablation_loss}
\input{tabs/ablation_model_size}

\para{Auxiliary losses.}
We finally investigate the contribution of each auxiliary loss in \cref{tab:ablation_loss}.
We observe a range of behaviors acros the auxiliary losses:
some significantly boost depth performance (\eg, 2D position, normal), others are critical for better camera pose estimation (\eg, confidence), and displacement has small improvements across the board.
